\documentclass[10pt,twocolumn]{article}
\usepackage[margin=0.75in,top=0.85in,bottom=0.85in]{geometry}
\usepackage[T1]{fontenc}
\usepackage[utf8]{inputenc}
\usepackage{microtype}
\usepackage{amsmath,amssymb}
\usepackage{booktabs,tabularx,array,multirow}
\usepackage{graphicx}
\usepackage{xcolor}
\usepackage[numbers,sort&compress]{natbib}
\usepackage{hyperref}
\usepackage{url}
\hypersetup{colorlinks=true,linkcolor=blue!50!black,citecolor=blue!50!black,urlcolor=blue!50!black}
\setlength{\parindent}{0.9em}
\setlength{\parskip}{0.2em}
\renewcommand{\arraystretch}{1.1}
\newcolumntype{C}{>{\centering\arraybackslash}p{1.6cm}}
\newcolumntype{Y}{>{\centering\arraybackslash}X}
\graphicspath{{../figures/}{./}}

\title{\textbf{Cross-Domain Limits of Hand-Crafted CoT-Surface Features}\\[4pt]
\large Strong Performance in Math, Narrower Gains in Science, Weak Results in Coding}
\author{Yuhan Chi\\Fudan University\\\texttt{yhchi25@m.fudan.edu.cn}}
\date{}

\begin{document}
\maketitle

\begin{abstract}
Can simple summaries of chain-of-thought (CoT) traces indicate whether a solution is correct? We study one family of such features---hand-crafted CoT-surface and token-trajectory descriptors---using \textbf{DeepSeek-R1-0528-Qwen3-8B} across math, science, and coding tasks.

Results differ substantially across domains. In math, the signal is strong: our probe achieves AUC-of-AUROC (AoA) 0.958 and AUROC@100\% 0.982, with dense-anchor curves plateauing early and best-of-$N=64$ reranking improving by $+10.0$ percentage points. Science shows a different pattern: the full probe reaches AoA 0.799 and AUROC@100\% 0.841 but does not exceed a simple token-confidence baseline (0.809 / 0.856), with permutation importance concentrated on confidence-based features.

Coding presents a distinct contrast. On LiveCodeBench-v5, the main probe achieves only AoA 0.434 and AUROC@100\% 0.407. Sweeping 83 coding-specific scalars peaks at AUROC 0.556. Grouped feature ablations remain near chance. A coding-specific CoT-only judge using 19 surface features achieves 5-fold GroupKFold out-of-fold AUC of $0.481 \pm 0.053$ and no best-of-$64$ gain ($-0.006$). Neither nonlinear models nor self-supervised pre-training substantially improve performance, and token-level de-knotting shows no benefit.

We therefore draw a narrower conclusion than earlier drafts: in this single-model study, hand-crafted CoT-surface features work well for math, partially for science, and poorly for coding correctness on unseen problems. This finding does not imply that raw-text classifiers, code-aware models, execution-aware selectors, or hidden-state probes must fail. Rather, it identifies one specific feature family whose semantics do not transfer cleanly to code verification.
\end{abstract}

\section{Introduction}

Process supervision and verifier-based reranking typically assume that answer quality leaves detectable traces in the generated chain of thought. This assumption often holds in mathematical reasoning, but its validity in coding remains unclear.

We examine a focused question: rather than asking whether all text-based verifiers fail on code, we investigate whether one particular family---cheap, interpretable CoT-surface features---transfers across domains. This family comprises token-confidence summaries, trajectory statistics, and activation-based descriptors that are computationally lightweight and straightforward to evaluate at test time.

Domain-specific results differ markedly. Mathematical reasoning exhibits strong early signals. Science problems show usable but narrower signals, primarily driven by token confidence. Coding presents a clear contrast: across multiple approaches---the main probe, extended scalar sweeps, grouped feature ablations, a coding-specific judge, nonlinear models, self-supervised pre-training, and token-level de-knotting---performance on unseen problems remains weak.

Our contribution is empirical: we demonstrate where this feature family excels (math), where it provides partial but limited utility (science), and where it fails to transfer (coding). This negative result has value---it shows that the chosen features do not reliably track correctness in all domains, even when they succeed in some.

\section{Related Work}

\paragraph{Process supervision and verifiers.}
Outcome verifiers and process reward models established that reasoning traces can be useful supervision signals \citep{cobbe2021training,lightman2023verify,wang2024mathshepherd}. ProcessBench evaluates step-level error detection in mathematical reasoning \citep{zheng2025processbench}, and generative verifier formulations cast reward modeling as next-token prediction \citep{zhang2025genrm}. Our paper asks a different question: before training a sophisticated verifier, what can be recovered from a very cheap feature family?

\paragraph{CoT faithfulness and self-correction.}
CoT traces are not always faithful explanations of model reasoning \citep{lanham2023faithfulness}. Related work on self-correction shows that reflective text does not reliably translate into better reasoning on its own \citep{huang2024selfcorrect,liu2024intrinsic,madaan2023selfrefine,shinn2023reflexion}. We add a cross-domain angle: even when self-corrective language is visible, the \emph{measurability} of correctness from that language depends on the task.

\paragraph{Code selection and execution-aware verification.}
Code-generation systems often rely on information channels stronger than CoT surface form. CodeT, for example, selects candidates by executing them against generated tests and comparing execution agreement \citep{chen2022codet}. Our setting is intentionally weaker: no execution feedback, no hidden-state probes, and no external API judge. The question is whether cheap hand-crafted CoT-surface features are already enough.

\paragraph{Measurement invariance.}
The psychometric literature on measurement invariance asks whether the same instrument measures the same latent construct across groups \citep{vandenberg2000review}. Probing work in NLP makes a related point: a probe can be informative without being semantically stable across settings \citep{belinkov2022probing}. We use that lens here. The same feature name can measure ``failure to converge'' in math and something much less useful in coding.

\section{Setup}
\label{sec:setup}

\subsection{Model, domains, and data}

All traces are generated by \textbf{DeepSeek-R1-0528-Qwen3-8B} at temperature 1.0.
\textbf{Math:} 7,680 aligned runs across AIME24, AIME25, BRUMO25, and HMMT25.
\textbf{Science:} 12,672 runs from GPQA.
\textbf{Coding:} 10,688 runs from LiveCodeBench-v5.

\subsection{Feature family}

The main probe uses five hand-crafted features computed from cached token streams and sparse activations.

\textit{Token-level summaries.}
\texttt{tok\_conf} is the mean negative log probability over visible prefix tokens. Four additional per-token arrays are available for auxiliary checks: \texttt{tok\_gini}, \texttt{tok\_logprob}, \texttt{tok\_neg\_entropy}, and \texttt{tok\_selfcert}.

\textit{Trajectory statistics.}
\texttt{traj\_continuity} is the mean Jaccard similarity between consecutive activation slices.
\texttt{traj\_novelty} is the mean $1 - \max_j \text{Jaccard}(s_i, s_j)$ over prior slices.
\texttt{traj\_reflection\_count} is the negated count of non-adjacent slice revisits above a fixed overlap threshold.

\textit{Neuron summary.}
\texttt{nc\_mean} is the mean row width of the sparse activation bank over the visible prefix.

\subsection{Probe pipeline and metrics}

The base probe is \texttt{StandardScaler} $\rightarrow$ \texttt{TruncatedSVD} $\rightarrow$ \texttt{LogisticRegression}, fit separately per domain. Evaluation uses problem-grouped splits.

The main early-stop summary is \textbf{AUC-of-AUROC (AoA)} over anchor positions $\{10,40,70,100\}$\%. We use AoA because the paper examines how signal emerges over the trace, not only at the end. To avoid over-relying on a single summary, we also report AUROC@100\% and use a denser 10-anchor check for robustness. Bootstrap 95\% confidence intervals on AoA use $B=10{,}000$ resamplings within each domain.

\section{Mathematical Reasoning: Strong Signal; Science: Narrower Signal}
\label{sec:positive}

Figure~\ref{fig:auroc} and Table~\ref{tab:positive} present the main cross-domain results. Math represents the clearest success case. Science shows positive but less clean results. Coding represents the failure case.

\begin{figure*}[t]
  \centering
  \includegraphics[width=\textwidth]{fig_auroc_by_anchor_v12_5.pdf}
  \caption{%
    \textbf{(a)} AUROC at each anchor position (10\%--100\% of trace visible) across domains.
    Mathematical reasoning rises early and stays high; science improves gradually; coding remains flat.
    \textbf{(b)} AoA with 95\% bootstrap CIs. The math--coding gap is substantial; science occupies an intermediate position.
  }
  \label{fig:auroc}
\end{figure*}

\begin{table}[t]
\centering
\footnotesize
\setlength{\tabcolsep}{3.0pt}
\begin{tabularx}{\columnwidth}{@{}lYYYY@{}}
\toprule
Domain & AoA & A100 & Fwd $\Delta$ & Bwd $\Delta$ \\
\midrule
Math    & \textbf{.958} [.931, .980] & \textbf{.982} & $-$.001 & $-$.006 \\
Science & \textbf{.799} [.775, .822] & \textbf{.841} & $-$.010 & $-$.001 \\
Coding  & .434 [.404, .464] & .407 & --- & --- \\
\midrule
\multicolumn{5}{l}{\textit{Baselines}} \\
\texttt{tok\_conf} & .506 / .809 / .506 & .490 / .856 / .490 & --- & --- \\
SC oracle$^\dagger$ & .953 / .921 / .961 & .953 / .921 / .961 & --- & --- \\
\bottomrule
\end{tabularx}
\caption{Main holdout results. AoA = AUC-of-AUROC; values in brackets are 95\% CIs from per-dataset bootstrap ($B=10{,}000$). A100 = AUROC at the 100\% anchor. Frozen-basis $\Delta$AUROC: Fwd = source-to-target transfer gap; Bwd = target-to-source. Baseline rows show math/science/coding values. $^\dagger$SC oracle = self-consistency upper bound using ground truth.}
\label{tab:positive}
\end{table}

\paragraph{Math: strong and structured signal.}
The math probe substantially outperforms the token-confidence baseline: AoA $0.958$ versus $0.685$, and AUROC@100\% $0.982$ versus $0.684$. This signal emerges early---on a denser 10-anchor grid, math reaches 95\% of its final AUROC by the 10\% anchor and plateaus by 50\%, achieving dense-anchor AoA $0.968$.

\paragraph{Science: usable but narrower than math.}
Science results stay above chance and improve with more context, yet differ qualitatively from math. The full probe achieves AoA $0.799$ and AUROC@100\% $0.841$, while a simple \texttt{tok\_conf} baseline reaches slightly higher values (AoA $0.809$, AUROC@100\% $0.856$). The science signal is genuine but appears largely confidence-driven rather than trajectory-structured.

\paragraph{Feature-importance contrast.}
Permutation importance highlights this difference. In math, \texttt{traj\_reflection\_count} dominates (importance $0.242$), with \texttt{traj\_continuity} ($0.092$) and \texttt{traj\_novelty} ($0.052$) following---a trajectory-driven pattern. In science, the leading features are \texttt{tok\_gini\_tail} ($0.038$), \texttt{tok\_conf\_recency} ($0.033$), \texttt{tok\_conf\_prefix} ($0.033$), and then \texttt{traj\_continuity} ($0.031$), reflecting a confidence-heavy rather than trajectory-structured signal.

\paragraph{Frozen-basis transfer.}
Cross-anchor transfer behaves differently across domains. Math exhibits near-lossless forward and backward transfer ($-0.001$ and $-0.006$), while science shows a larger forward gap ($-0.010$). This suggests math has a more stable quality geometry across trace positions.

\paragraph{Dense-anchor robustness.}
The four-anchor summary captures the main pattern. On a denser 10-anchor grid, the domain ordering remains unchanged: AoA reaches $0.968$ for math, $0.827$ for science, and $0.432$ for coding. Science continues improving into later anchors, whereas coding stays weak throughout.

\section{Coding: Persistent Weakness Across CoT-Surface Variants}
\label{sec:coding}

For coding, these hand-crafted CoT-surface features fail to predict correctness on unseen problems. We demonstrate this through multiple grouped cross-validation evaluations.

\subsection{Main probe and scalar screen}

The main coding probe achieves AoA $0.434$ and AUROC@100\% $0.407$, below the \texttt{tok\_conf} baseline (AoA $0.506$, AUROC@100\% $0.490$). Screening 83 coding-specific scalars yields limited improvement: the best single feature, \texttt{struct\_density}, reaches oriented AUROC $0.565$, followed by \texttt{cf\_n\_ifs} ($0.563$) and \texttt{cot\_trace\_count} ($0.558$). Even the optimal top-4 combination achieves only AUROC $0.556$.

\paragraph{Interpretation.}
These features are not pure noise---they exhibit weak but detectable correlations. However, the strongest effects remain small and semantically mixed. For instance, deeper branching and more \texttt{if}-statements weakly correlate positively with correctness, while denser structural narration shows weak negative correlation. While variation exists, it does not reliably separate correct from incorrect solutions.

\subsection{Grouped family ablations}

To test whether failure stems from a specific subset rather than the entire family, we evaluate grouped feature subsets using 5-fold GroupKFold on the 10,688 coding runs. No subset performs well: \texttt{traj\_only} achieves AUROC $0.509$, \texttt{token\_plus\_traj} reaches $0.501$, and the full 30-feature surface family tops out at $0.506$ in its best configuration. All remain in weak discrimination regimes.

\subsection{Coding-specific CoT-only run judge}

We then construct a coding-specific CoT-only judge using 19 surface features: planning and execution density, vocabulary overlap, structural flags, length summaries, one neuron count, and four one-hot \texttt{cot\_mode} indicators. This model remains cheap and interpretable, evaluated on unseen \texttt{problem\_id}s using 5-fold GroupKFold.

Performance remains weak. Mean out-of-fold AUC is $0.481 \pm 0.053$. By difficulty band, results stay near chance: hard $0.515$, mid $0.470$, easy $0.509$. In best-of-$64$ reranking, the judge achieves pass@1 $0.581$ against a random baseline of $0.587$ ($\Delta=-0.006$). Surface cues that appear meaningful within problems do not generalize to unseen problems.

\subsection{Nonlinear and self-supervised checks}

A one-hidden-layer MLP on the 25-dimensional feature family produces similar results: mean-anchor AUROC $0.499$ and AUROC@100\% $0.507$.

Self-supervised pre-training provides further evidence. Cross-domain SSL v1 collapses at one rank without meaningful coding improvement. Domain-specific SSL v2 avoids collapse and offers modest low-label regularization, yet the full-label coding ceiling remains weak: the best v2 configuration achieves AUROC $0.454$ at label fraction 1.0, peaking at $0.537$ when label fraction is 0.05. We interpret SSL as indicating that these objectives do not uncover additional label-relevant structure in this feature space, not as proof that no useful structure exists.

\begin{figure*}[t]
  \centering
  \includegraphics[width=\textwidth]{fig_ssl_ceiling_v12_5.pdf}
  \caption{%
    \textbf{SSL ceiling in coding vs.\ science.}
    \textbf{(a)} In coding, the full-label ceiling remains weak across SSL variants.
    \textbf{(b)} Science: SSL provides greater low-label benefits, consistent with a domain where features carry label-relevant structure.
  }
  \label{fig:ssl}
\end{figure*}

\subsection{De-knotting remains neutral in coding}

One alternative explanation posits that coding traces contain noisy circular spans that obscure an otherwise useful signal. We test this by removing GLM-labeled knot spans from all cached per-token arrays, not just from the text.

\begin{figure*}[t]
  \centering
  \includegraphics[width=\textwidth]{fig_deknot_alldomains_v12_5.pdf}
  \caption{%
    \textbf{De-knotting effects across domains.}
    Removing knot tokens degrades math performance, shows minimal effect in science, and remains neutral for coding.
    The coding weakness is therefore not attributable to removable surface clutter.
  }
  \label{fig:deknot}
\end{figure*}

Results differ by domain. In math, removing knot tokens degrades AUROC ($0.502 \rightarrow 0.453$, $\Delta=-0.049$), indicating these spans contribute to the signal. Science shows minimal change ($\Delta=-0.006$). Coding remains neutral ($0.460 \rightarrow 0.466$, $\Delta=+0.006$). De-knotting does not rescue coding performance.

\begin{table}[t]
\centering
\footnotesize
\setlength{\tabcolsep}{3.0pt}
\begin{tabularx}{\columnwidth}{@{}>{\raggedright\arraybackslash}X>{\raggedright\arraybackslash}p{2.3cm}>{\raggedright\arraybackslash}p{1.45cm}@{}}
\toprule
Check & Best coding result & Reading \\
\midrule
Main probe & AoA .434 / A100 .407 & Below \texttt{tok\_conf} \\
83-scalar screen & AUROC .556 & Weak proxy only \\
Family ablations & AUROC .509 (\texttt{traj\_only}) & No stable subset \\
CoT-only run judge & OOF AUC .481 $\pm$ .053 & Fails on unseen problems \\
MLP & AUROC .499--.507 & No meaningful gain \\
SSL v2 & .537 at lf=.05; .454 at lf=1.0 & Low-label regularizer only \\
De-knotting & $\Delta = +.006$ & Neutral \\
\bottomrule
\end{tabularx}
\caption{Coding-side robustness checks. Results are drawn from grouped or held-out evaluation in the artifact tables. No approach produces a strong, generalizable coding verifier using only cheap CoT-surface features.}
\label{tab:convergence}
\end{table}

\section{Downstream Impact: Best-of-\texorpdfstring{$N$}{N} Reranking}
\label{sec:reranking}

The cross-domain probe results affect more than just metrics---they change downstream selection behavior.

\begin{figure}[t]
  \centering
  \includegraphics[width=\linewidth]{fig_reranking_v12_5.pdf}
  \caption{%
    Best-of-$N=64$ probe reranking results. Math gains $+10.0$pp, science gains $+8.0$pp,
    and coding shows no meaningful improvement.
  }
  \label{fig:reranking}
\end{figure}

\begin{table}[t]
\centering
\footnotesize
\setlength{\tabcolsep}{3.5pt}
\begin{tabularx}{\columnwidth}{@{}l c Y Y c@{}}
\toprule
Domain & $n$ & Random & Probe & $\Delta$ \\
\midrule
Math    & 120 & .642 [.572, .712] & .742 [.658, .817] & \textbf{$+$.100} \\
Science & 198 & .602 [.550, .652] & .682 [.616, .748] & \textbf{$+$.080} \\
Coding  & 167 & .617 [.545, .689] & .611 [.539, .683] & $-$.006 \\
\bottomrule
\end{tabularx}
\caption{Best-of-$N=64$ reranking pass@1. Values in brackets are 95\% bootstrap CIs from problem-level resampling.}
\label{tab:reranking}
\end{table}

Math and science show substantial gains; coding does not. The coding-specific CoT-only judge performs similarly: out-of-fold reranking yields $0.581$ against the same $0.587$ random baseline. For this feature family, weak run-level discrimination translates to no useful selection improvement.

\section{Discussion}
\label{sec:why}

\paragraph{What we claim.}
Within this single-model study, hand-crafted CoT-surface features are domain-specific tools: they work well for math, partially for science, and poorly for coding.

\paragraph{What we do not claim.}
We do not argue that correctness is fundamentally undetectable in text, nor that raw-text classifiers, code-aware models, execution-aware selectors, hidden-state probes, or LLM judges must fail. A more plausible explanation is that coding requires different signals, potentially execution-based ones.

\paragraph{Why math and coding diverge.}
In math, visible hesitation and convergence patterns often reflect underlying reasoning quality. Reflection loops, continuity breaks, and recency-weighted confidence track whether the reasoning path has stabilized. In coding, correctness depends on program execution rather than narrative coherence. A clean explanation can implement incorrect logic, while a messy explanation may produce correct code. The text trace is thus further removed from the target variable.

\paragraph{Why science occupies the middle ground.}
Science problems share some verbal uncertainty cues with math, but our results suggest the usable signal comes primarily from token confidence rather than strong trajectory structure. This explains why science appears positive in aggregate but weaker on closer inspection.

\paragraph{Future work.}
Subsequent work should test a second model, include raw-text or LLM-judge baselines, and add an execution-aware positive control. These experiments would clarify whether the current failure is specific to this feature family, this model, or both.

\section{Limitations and Conclusion}
\label{sec:limits}

\paragraph{Limitations.}
First, all experiments use a single model (\textbf{DeepSeek-R1-0528-Qwen3-8B}). Second, our strongest coding check remains a hand-crafted CoT-only judge; we do not include raw-text classifiers or external LLM judges. Third, the coding experiments lack an execution-aware positive control. Fourth, our SSL conclusion is objective-specific: it shows failure under the tested pretraining losses, not a universal impossibility.

\paragraph{Conclusion.}
We draw a narrower conclusion than earlier drafts. The evidence supports this statement:

\begin{quote}
\textbf{Hand-crafted CoT-surface features transfer well to math, partly to science, and poorly to coding correctness on unseen problems in this single-model study.}
\end{quote}

This negative result remains valuable: it shows that cheap CoT measurements should be treated as domain-specific instruments, not general-purpose correctness proxies. For coding, at least with this model and feature family, they are the wrong tool.

\paragraph{Code and artifacts.}
The public repository for this project is \url{https://github.com/Chi-Shan0707/code-not-text}. In the local snapshot used for this draft, code and paper assets remain under \texttt{workshop/cotknot/}. Paper sources are in \texttt{workshop/cotknot/paper/}. Quantitative tables used here come from \texttt{results/tables/} and \texttt{workshop/cotknot/results/tables/}. The coding-side robustness checks in this revision rely in particular on four artifacts: \texttt{results/tables/coding\_feature\_family\_ablation.csv}, \texttt{results/tables/cot\_run\_judge\_scores.csv}, \texttt{results/tables/cot\_run\_judge\_rerank.csv}, and \texttt{results/models/cot\_run\_judge.summary.json}.

\bibliographystyle{plainnat}
\begin{thebibliography}{16}

\bibitem{wei2022cot}
Jason Wei et al.
\newblock Chain-of-thought prompting elicits reasoning in large language models.
\newblock In \emph{NeurIPS}, 2022.

\bibitem{cobbe2021training}
Karl Cobbe et al.
\newblock Training verifiers to solve math word problems.
\newblock \emph{arXiv:2110.14168}, 2021.

\bibitem{lightman2023verify}
Hunter Lightman et al.
\newblock Let's verify step by step.
\newblock \emph{arXiv:2305.20050}, 2023.

\bibitem{wang2024mathshepherd}
Peiyi Wang et al.
\newblock Math-Shepherd: Verify and reinforce {LLM}s step-by-step without human annotations.
\newblock In \emph{ACL}, 2024.

\bibitem{zheng2025processbench}
Chujie Zheng et al.
\newblock ProcessBench: Identifying process errors in mathematical reasoning.
\newblock In \emph{ACL}, 2025.

\bibitem{zhang2025genrm}
Lunjun Zhang et al.
\newblock Generative verifiers: Reward modeling as next-token prediction.
\newblock In \emph{ICLR}, 2025.

\bibitem{lanham2023faithfulness}
Tamera Lanham et al.
\newblock Measuring faithfulness in chain-of-thought reasoning.
\newblock \emph{arXiv:2307.13702}, 2023.

\bibitem{huang2024selfcorrect}
Jie Huang et al.
\newblock Large language models cannot self-correct reasoning yet.
\newblock \emph{arXiv:2310.01798}, 2023.

\bibitem{liu2024intrinsic}
Dancheng Liu et al.
\newblock Large language models have intrinsic self-correction ability.
\newblock In \emph{NeurIPS}, 2024.

\bibitem{madaan2023selfrefine}
Aman Madaan et al.
\newblock Self-Refine: Iterative refinement with self-feedback.
\newblock \emph{arXiv:2303.17651}, 2023.

\bibitem{shinn2023reflexion}
Noah Shinn et al.
\newblock Reflexion: Language agents with verbal reinforcement learning.
\newblock \emph{arXiv:2303.11366}, 2023.

\bibitem{chen2022codet}
Bei Chen et al.
\newblock CodeT: Code generation with generated tests.
\newblock \emph{arXiv:2207.10397}, 2022.

\bibitem{vandenberg2000review}
Robert E.\ Vandenberg and Charles E.\ Lance.
\newblock A review and synthesis of the measurement invariance literature.
\newblock \emph{Organizational Research Methods}, 3(1):4--70, 2000.

\bibitem{belinkov2022probing}
Yonatan Belinkov.
\newblock Probing classifiers: Promises, shortcomings, and advances.
\newblock \emph{Computational Linguistics}, 48(1):207--219, 2022.

\end{thebibliography}

\end{document}
